\subsection{Spécification de la fonction d'utilité}
\label{sec:spec-utilite}

Cette sous-section décrit la fonction d'utilité telle qu'elle est implémentée dans
\texttt{src/logminer/agents/intelligent\_runtime.py}. Soit un agent $a$ et une tâche
$\tau$ de type $t$. On note $\mathcal{K}_a(\tau)$ l'ensemble des capacités de $a$ qui
supportent $t$, $n_a$ son nombre de tâches actives,
$P_a=\max(1,\texttt{max\_parallel\_tasks})$ et
$\rho_a=\min(1,n_a/P_a)$ son taux de charge. Pour le type $t$, $s_t$ et $e_t$
désignent les nombres cumulés d'exécutions terminées respectivement avec le statut
\texttt{ok} et avec un autre statut ; $D_t$ contient les durées des cinquante dernières
exécutions. Enfin, $\mu\in\{0,1\}$ indique si la mémoire est activée.

\paragraph{Règles de refus.}
Avant de calculer une offre, l'agent vérifie successivement la présence d'une capacité
compatible, des modèles et dépendances exigés, son état de santé et sa capacité de charge.
Ces contrôles produisent respectivement les motifs
\texttt{capability\_missing}, \texttt{model\_missing},
\texttt{dependency\_unavailable}, \texttt{agent\_unhealthy} et
\texttt{agent\_overloaded}. Une sixième règle produit
\texttt{low\_expected\_utility} lorsque $U<u_{\min}$, avec
$u_{\min}=\min(1,\max(0,\texttt{minimum\_utility}))$. Sa valeur par défaut est 0 :
elle ne filtre donc aucune offre tant qu'un seuil strictement positif n'est pas configuré.

\paragraph{Composantes.}
Lorsque les cinq premiers contrôles sont satisfaits, les composantes sont calculées comme
indiqué au tableau~\ref{tab:spec-utilite}.

\begin{table}[H]
\centering
\caption{Définition opérationnelle des composantes de l'utilité}
\label{tab:spec-utilite}
\small
\begin{tabularx}{\textwidth}{p{0.09\textwidth} p{0.55\textwidth} X}
\toprule
\textbf{Comp.} & \textbf{Calcul implémenté} & \textbf{Interprétation} \\
\midrule
$C$ & $\min\!\left(1,\max\!\left(0,\max_{k\in\mathcal{K}_a(\tau)}c_k\right)\right)$
& Confiance maximale parmi les capacités compatibles. \\
$H$ & $\mu\dfrac{s_t+1}{s_t+e_t+2}+(1-\mu)\dfrac{1}{2}$
& Fiabilité historique lissée ; valeur neutre $0{,}5$ sans mémoire. \\
$M$ & $1$
& Compatibilité du modèle ; elle agit comme filtre préalable et non comme score variable. \\
$A$ & $\max(0,1-\rho_a)$
& Disponibilité dérivée de la capacité d'exécution restante. \\
$L$ & $\max(0,1-\rho_a)$
& Charge, calculée par la même expression que $A$. \\
$T$ & $1/(1+\bar d_t)$ si $\mu=1$ et $D_t\neq\varnothing$, $0{,}5$ sinon
& Latence historique, avec $\bar d_t$ moyenne des durées en secondes. \\
\bottomrule
\end{tabularx}
\end{table}

La compatibilité des modèles n'intervient ainsi que comme filtre : $M=1$ pour toute offre
émise. De même, $A$ et $L$ ne se distinguent dans le calcul que par leur poids. Enfin,
$T$ dépend de l'unité de temps retenue, la seconde, puisque la durée moyenne n'est pas
rapportée à une durée de référence.

\paragraph{Agrégation.}
L'utilité est la combinaison convexe
\[
U=\min\!\left(1,\max\!\left(0,
w_CC+w_HH+w_MM+w_AA+w_LL+w_TT\right)\right),
\]
avec, par défaut,
$(w_C,w_H,w_M,w_A,w_L,w_T)=(0{,}30;0{,}20;0{,}20;0{,}15;0{,}10;0{,}05)$.
Les poids sont configurables, positifs ou nuls, et leur somme doit être égale à 1 à
$10^{-9}$ près. Le choix des six composantes et de leurs poids est arbitraire : il a été
fixé a priori par l'auteur, selon ce qui lui paraissait cohérent pour l'application (la
capacité déclarée en premier, puis la fiabilité observée et la compatibilité du modèle, la
disponibilité et la charge ensuite, la latence en dernier). Il n'est ni dérivé d'un modèle,
ni optimisé, ni comparé à d'autres pondérations ; la redondance de $A$ et de $L$ et la
constance de $M$ en résultent. Comme $M=1$ et $A=L$, toute offre émise vérifie, avec ces poids,
\[
U=0{,}20+0{,}30C+0{,}20H+0{,}25(1-\rho_a)+0{,}05T.
\]
Dans les campagnes Redis, $P_a=1$. Une offre suppose alors $n_a=0$, donc $\rho_a=0$ et
\[
U=0{,}45+0{,}30C+0{,}20H+0{,}05T.
\]
Sans mémoire, $H=T=0{,}5$ et l'utilité se réduit à $U=0{,}575+0{,}30C$ : elle reste
strictement inférieure à 1 et ne dépend que de la confiance déclarée. Une confiance de
$0{,}9$ donne ainsi $U=0{,}845$, et les confiances $0{,}95$ et $0{,}78$ donnent
$0{,}86$ et $0{,}809$.
Le paramètre \texttt{minimum\_utility} n'y est pas transmis ; le refus pour faible utilité
y reste donc inactif. L'utilité est ainsi un classement à poids fixes : dans les campagnes
Redis, trois de ses six termes sont constants ($M=1$, et $A=L=1$ puisque $\rho_a=0$), et
seuls $C$, $H$ et $T$ différencient les offres, $H$ et $T$ uniquement lorsque la mémoire est
activée. La décision des agents ne relève donc pas d'un raisonnement, mais d'une règle de
classement explicite et vérifiable.

\paragraph{Mise à jour et rôle décisionnel.}
$H$ et $T$ évoluent après une exécution locale par l'appel à \texttt{learn}, uniquement
si la mémoire est activée. Les messages \texttt{FEEDBACK} et \texttt{REJECT} ne modifient
pas cette mémoire. Dans le Contract Net local comme dans son transport Redis, $U$ sert à
classer les offres par ordre décroissant ; une égalité est départagée par l'identifiant de
l'agent. Le coordinateur ne compare pas $U$ à un seuil supplémentaire.

Pour la poursuite locale d'un objectif, l'utilité n'est qu'un terme du score
\begin{multline*}
S=\pi+20U-100\,\mathbf{1}[\tau\in\mathcal{E}]
-1{,}5r\,\mathbf{1}[r\geq4] \\
{}+\bigl(10+\max(0,5-\delta)\bigr)\mathbf{1}[\delta<5]
-2\min(\alpha,5)\mathbf{1}[\alpha>0],
\end{multline*}
où $\pi$ est la priorité, $\mathcal{E}$ l'historique récent, $r$ la répétition du type,
$\delta$ l'échéance et $\alpha$ le nombre de tentatives. La priorité et les correctifs
peuvent donc dominer $20U$. L'autonomie demeure bornée par défaut à 20 actions, une
profondeur de 5, deux tentatives par tâche et 100 tâches en attente. À épuisement de ces
limites, l'objectif est bloqué, échoue ou passe à l'état
\texttt{budget\_exhausted} selon le garde-fou concerné.
